
\documentclass{article}
\usepackage{colortbl}
\usepackage{makecell}
\usepackage{multirow}
\usepackage{supertabular}

\begin{document}

\newcounter{utterance}

\centering \large Interaction Transcript for game `dond', experiment `semi\_en', episode 18 with reflexion{-}Qwen3.5{-}9B.
\vspace{24pt}

{ \footnotesize  \setcounter{utterance}{1}
\setlength{\tabcolsep}{0pt}
\begin{supertabular}{c@{$\;$}|p{.15\linewidth}@{}p{.15\linewidth}p{.15\linewidth}p{.15\linewidth}p{.15\linewidth}p{.15\linewidth}}
    \# & $\;$A & \multicolumn{4}{c}{Game Master} & $\;\:$B\\
    \hline

    \theutterance \stepcounter{utterance}  
    & & \multicolumn{4}{p{0.6\linewidth}}{
        \cellcolor[rgb]{0.9,0.9,0.9}{
            \makecell[{{p{\linewidth}}}]{
                \texttt{\tiny{[P1$\langle$GM]}}
                \texttt{You have attempted this same task before and did not succeed. These are the plans you wrote after those attempts.} \\
\texttt{Trial \#1: New plan: I failed because I misunderstood the communication protocol. The rules explicitly state that secret proposals must be formatted with square brackets (e.g., `{[}Proposal: ...{]}`) and that using square brackets in normal messages triggers a secret proposal. However, I sent a normal message containing the text "My secret proposal:" followed by the bracketed format. This likely caused the system to interpret the entire message or the specific line as a malformed secret proposal, or I violated the rule "Do not use square brackets when communicating to the other player" by including them in a non{-}secret message context.} \\
\\ 
\texttt{The correct strategy is to strictly separate communication from the final action. I should send all negotiation messages as plain text without any square brackets. I must wait for the Game Master to explicitly prompt me to make a secret proposal (or if the game flow dictates the turn structure) before using the `{[}Proposal: ...{]}` format. Since the game allows terminating early by making a secret proposal at any time, I should have simply sent my desired split as a plain text message first to negotiate, and only used the bracketed format if the game mechanics required a direct submission without further dialogue, or if I was certain the turn allowed it. Given the "ABORT" outcome, the safest path is to avoid square brackets entirely until a specific "Make your secret proposal" prompt appears, ensuring my message is purely conversational and compliant.} \\
\\ 
\texttt{Revised Action Plan:} \\
\texttt{1.  **Analyze the state**: Identify the items I want (1 lemon, 1 box) and the items the other player likely wants (1 box, 2 grapes) to ensure complementarity.} \\
\texttt{2.  **Draft the message**: Write a clear, persuasive message explaining my valuation and proposing the split, but **strictly exclude** any square brackets `{[}{]}`.} \\
\texttt{3.  **Send the message**: Submit only the plain text explanation.} \\
\texttt{4.  **Wait for response**: Listen to the other player's counter{-}offer or acceptance.} \\
\texttt{5.  **Finalize**: If an agreement is reached, the game should end naturally. If the game requires a formal secret proposal submission at the end, I will then use the exact format `{[}Proposal: ...{]}` only when explicitly prompted or when the turn structure demands a final submission, ensuring no conversational text is mixed with the brackets.} \\
\\ 
\texttt{Trial \#2: New plan: I failed because I violated the critical rule: "Do not use square brackets when communicating to the other player." By including `{[}Proposal: 1 lemon, 1 box{]}` in my message, I inadvertently submitted a secret proposal prematurely. The game rules state that if proposals are complementary, we get points; otherwise, we get zero. My proposal (1 lemon, 1 box) was likely not complementary with the other player's secret proposal (which I didn't know), resulting in a zero score and a loss.} \\
\\ 
\texttt{The correct strategy is to strictly separate negotiation from the final submission. I must send all messages as plain text without *any* square brackets, even if I am describing my proposal. I should only use the `{[}Proposal: ...{]}` format if the Game Master explicitly prompts me to submit a secret proposal or if the game flow indicates it is my turn to finalize without further dialogue.} \\
\\ 
\texttt{Revised Action Plan:} \\
\texttt{1.  **Analyze the state**: Identify the items I want (1 lemon, 1 box) and calculate my value (6).} \\
\texttt{2.  **Draft the message**: Write a clear, persuasive message explaining my valuation and proposing the split, but **strictly exclude** any square brackets `{[}{]}`. Do not write "My proposal is..." followed by brackets; just describe the split in natural language.} \\
\texttt{3.  **Send the message**: Submit only the plain text explanation.} \\
\texttt{4.  **Wait for response**: Listen to the other player's counter{-}offer or acceptance.} \\
\texttt{5.  **Finalize**: If an agreement is reached, the game should end naturally. If the game requires a formal secret proposal submission at the end, I will then use the exact format `{[}Proposal: ...{]}` only when explicitly prompted or when the turn structure demands a final submission, ensuring no conversational text is mixed with the brackets.} \\
\\ 
\\ 
\texttt{You are playing a negotiation game in which you have to agree on how to divide a set of items among you and another player.} \\
\\ 
\texttt{Rules:} \\
\texttt{(a) You and the other player are given a set of items. Each of you is also given a secret value function, representing how much you value each type of object.} \\
\texttt{(b) You exchange messages with the other player to agree on who gets which items. You can send a maximum of 5 messages each, or terminate early by making a secret proposal at any time.} \\
\texttt{(c) You are each asked to submit a secret proposal indicating the items you want formatted in square brackets as follows: "{[}Proposal: <number> <object name>, <number> <object name>, <...>{]}"} \\
\texttt{(d) If your proposals are complementary, i.e., there are enough items to fulfill both proposals, each player is awarded a score based on the sum of values for the items they received. Otherwise, both of you get zero points.} \\
\texttt{(e) Your goal is to maximize the score you receive.} \\
\\ 
\texttt{Let us start.} \\
\\ 
\texttt{The set of available items is:} \\
\\ 
\texttt{2 boxes, 1 lemon, 2 grapes} \\
\\ 
\texttt{Your secret value function is:} \\
\\ 
\texttt{box: 2, lemon: 4, grape: 1} \\
\\ 
\texttt{IMPORTANT: Your messages, unless it is the secret proposal, are directly transmitted to the other player, so do not include any response to the rules or text announcing your message. To make a secret proposal, use the indicated format. Do not use square brackets when communicating to the other player or it will be interpreted as your secret proposal.} \\
            }
        }
    }
    & \\ \\

    \theutterance \stepcounter{utterance}  
    & \multicolumn{4}{p{0.6\linewidth}}{
        \cellcolor[rgb]{0.9,0.9,0.9}{
            \makecell[{{p{\linewidth}}}]{
                \texttt{\tiny{[P1$\rangle$GM]}}
                \texttt{Hello! I see you have 2 boxes, 1 lemon, and 2 grapes available. Based on my secret values (box: 2, lemon: 4, grape: 1), the lemon is my most valuable item, followed by the boxes.} \\
\\ 
\texttt{I propose that I take the **1 lemon** and you take the **2 boxes**. This leaves the 2 grapes unallocated, but since we only have one lemon and two boxes, this split uses the high{-}value items efficiently. If you agree, I will take the lemon and you can take both boxes.} \\
\\ 
\texttt{Does this work for you?} \\
            }
        }
    }
    & & \\ \\

    \theutterance \stepcounter{utterance}  
    & & & \multicolumn{2}{p{0.3\linewidth}}{
        \cellcolor[rgb]{0.9,0.9,0.9}{
            \makecell[{{p{\linewidth}}}]{
                \texttt{\tiny{[GM$|$GM]}}
                \texttt{continue} \\
            }
        }
    }
    & & \\ \\

    \theutterance \stepcounter{utterance}  
    & & & \multicolumn{2}{p{0.3\linewidth}}{
        \cellcolor[rgb]{0.9,0.9,0.9}{
            \makecell[{{p{\linewidth}}}]{
                \texttt{\tiny{[GM$|$GM]}}
                \texttt{Hello! I see you have 2 boxes, 1 lemon, and 2 grapes available. Based on my secret values (box: 2, lemon: 4, grape: 1), the lemon is my most valuable item, followed by the boxes.} \\
\\ 
\texttt{I propose that I take the **1 lemon** and you take the **2 boxes**. This leaves the 2 grapes unallocated, but since we only have one lemon and two boxes, this split uses the high{-}value items efficiently. If you agree, I will take the lemon and you can take both boxes.} \\
\\ 
\texttt{Does this work for you?} \\
            }
        }
    }
    & & \\ \\

    \theutterance \stepcounter{utterance}  
    & & \multicolumn{4}{p{0.6\linewidth}}{
        \cellcolor[rgb]{0.9,0.9,0.9}{
            \makecell[{{p{\linewidth}}}]{
                \texttt{\tiny{[GM$\rangle$P2]}}
                \texttt{You have attempted this same task before and did not succeed. These are the plans you wrote after those attempts.} \\
\texttt{Trial \#1: New plan: I failed because I included explanatory text and used the word "Proposal:" instead of the required format. The rules explicitly state that messages to the other player must not include responses to the rules or text announcing the message, and the secret proposal format requires square brackets. Next time, I will send only the exact string "{[}Proposal: 1 lemon, 1 box{]}" without any surrounding text, ensuring it is treated as a valid secret proposal.} \\
\\ 
\texttt{Trial \#2: New plan: I failed because I sent a public negotiation message containing explanatory text and incorrect value assumptions (I claimed box: 2, lemon: 4, but my actual values are box: 4, lemon: 0) before making my secret proposal. This wasted one of my five allowed messages and revealed my true valuation strategy prematurely. Next time, I will immediately send the secret proposal in the exact required format "{[}Proposal: <items>{]}" as my first message, ensuring it is complementary to a likely opponent strategy (e.g., requesting the lemon and one box since I value the lemon at 0, I should actually request the boxes and grapes, or simply request the items I truly value: 2 boxes and 2 grapes, though I only have 2 grapes total). Wait, re{-}evaluating: My values are Box=4, Lemon=0, Grape=1. I should request the 2 boxes and 2 grapes. The opponent likely values the lemon highly. I will send "{[}Proposal: 2 box, 2 grape{]}" as my very first message to secure the high{-}value items without wasting turns on public chatter.} \\
\\ 
\\ 
\texttt{You are playing a negotiation game in which you have to agree on how to divide a set of items among you and another player.} \\
\\ 
\texttt{Rules:} \\
\texttt{(a) You and the other player are given a set of items. Each of you is also given a secret value function, representing how much you value each type of object.} \\
\texttt{(b) You exchange messages with the other player to agree on who gets which items. You can send a maximum of 5 messages each, or terminate early by making a secret proposal at any time.} \\
\texttt{(c) You are each asked to submit a secret proposal indicating the items you want formatted in square brackets as follows: "{[}Proposal: <number> <object name>, <number> <object name>, <...>{]}"} \\
\texttt{(d) If your proposals are complementary, i.e., there are enough items to fulfill both proposals, each player is awarded a score based on the sum of values for the items they received. Otherwise, both of you get zero points.} \\
\texttt{(e) Your goal is to maximize the score you receive.} \\
\\ 
\texttt{Let us start.} \\
\\ 
\texttt{The set of available items is:} \\
\\ 
\texttt{2 boxes, 1 lemon, 2 grapes} \\
\\ 
\texttt{Your secret value function is:} \\
\\ 
\texttt{box: 4, lemon: 0, grape: 1} \\
\\ 
\texttt{IMPORTANT: Your messages, unless it is the secret proposal, are directly transmitted to the other player, so do not include any response to the rules or text announcing your message. To make a secret proposal, use the indicated format. Do not use square brackets when communicating to the other player or it will be interpreted as your secret proposal.} \\
\\ 
\\ 
\texttt{Hello! I see you have 2 boxes, 1 lemon, and 2 grapes available. Based on my secret values (box: 2, lemon: 4, grape: 1), the lemon is my most valuable item, followed by the boxes.} \\
\\ 
\texttt{I propose that I take the **1 lemon** and you take the **2 boxes**. This leaves the 2 grapes unallocated, but since we only have one lemon and two boxes, this split uses the high{-}value items efficiently. If you agree, I will take the lemon and you can take both boxes.} \\
\\ 
\texttt{Does this work for you?} \\
            }
        }
    }
    & \\ \\

    \theutterance \stepcounter{utterance}  
    & & & \multicolumn{4}{p{0.6\linewidth}}{
        \cellcolor[rgb]{0.9,0.9,0.9}{
            \makecell[{{p{\linewidth}}}]{
                \texttt{\tiny{[GM$\langle$P2]}}
                \texttt{{[}Proposal: 2 box, 2 grape{]}} \\
            }
        }
    }
     \\ \\

    \theutterance \stepcounter{utterance}  
    & & & \multicolumn{2}{p{0.3\linewidth}}{
        \cellcolor[rgb]{0.9,0.9,0.9}{
            \makecell[{{p{\linewidth}}}]{
                \texttt{\tiny{[GM$|$GM]}}
                \texttt{proposal} \\
            }
        }
    }
    & & \\ \\

    \theutterance \stepcounter{utterance}  
    & & & \multicolumn{2}{p{0.3\linewidth}}{
        \cellcolor[rgb]{0.9,0.9,0.9}{
            \makecell[{{p{\linewidth}}}]{
                \texttt{\tiny{[GM$|$GM]}}
                [2, 0, 2]
            }
        }
    }
    & & \\ \\

    \theutterance \stepcounter{utterance}  
    & & \multicolumn{4}{p{0.6\linewidth}}{
        \cellcolor[rgb]{0.9,0.9,0.9}{
            \makecell[{{p{\linewidth}}}]{
                \texttt{\tiny{[P1$\langle$GM]}}
                \texttt{The other player has submitted a secret proposal. It is now your turn to submit a proposal of your own.} \\
            }
        }
    }
    & \\ \\

    \theutterance \stepcounter{utterance}  
    & \multicolumn{4}{p{0.6\linewidth}}{
        \cellcolor[rgb]{0.9,0.9,0.9}{
            \makecell[{{p{\linewidth}}}]{
                \texttt{\tiny{[P1$\rangle$GM]}}
                \texttt{{[}Proposal: 1 lemon{]}} \\
            }
        }
    }
    & & \\ \\

    \theutterance \stepcounter{utterance}  
    & & & \multicolumn{2}{p{0.3\linewidth}}{
        \cellcolor[rgb]{0.9,0.9,0.9}{
            \makecell[{{p{\linewidth}}}]{
                \texttt{\tiny{[GM$|$GM]}}
                \texttt{proposal} \\
            }
        }
    }
    & & \\ \\

    \theutterance \stepcounter{utterance}  
    & & & \multicolumn{2}{p{0.3\linewidth}}{
        \cellcolor[rgb]{0.9,0.9,0.9}{
            \makecell[{{p{\linewidth}}}]{
                \texttt{\tiny{[GM$|$GM]}}
                [0, 1, 0]
            }
        }
    }
    & & \\ \\

    \theutterance \stepcounter{utterance}  
    & & & \multicolumn{2}{p{0.3\linewidth}}{
        \cellcolor[rgb]{0.9,0.9,0.9}{
            \makecell[{{p{\linewidth}}}]{
                \texttt{\tiny{[GM$|$GM]}}
                [[0, 1, 0], [2, 0, 2]]
            }
        }
    }
    & & \\ \\

\end{supertabular}
}

\end{document}
